\input{Iteration-V2/Chapter6/chapter-06-11}

\subsection{Revision 30 temporal-alert limits and next validation}

Temporal co-occurrence is necessarily an incomplete proxy. A combined location, sensor, or transfer pattern may be benign, expected by the user, technically misclassified, partially observed, or governed by a lawful purpose not present in the local evidence. Conversely, an actual high-risk flow can evade this prototype when it occurs outside a profile's window, uses an unsupported interface, is hidden by platform restrictions, or lacks observable bridge evidence. The bounded ledger also intentionally trades historical completeness for local resource control. These are reasons to preserve uncertainty and notification-only operation, not reasons to silently broaden the rule.

Future regulation packs should be authored through controlled legal and engineering review, with versioned source material, traceable interpretation notes, test fixtures for both matches and non-matches, and a staged rollback path. A rule author must set the window and required combination for the assessment item and document why those parameters are appropriate; the generic engine should not invent a default. Pack loading must continue to validate the converter boundary so that unsupported observation types, empty legal references, unsafe counts, or implausible windows cannot activate accidental alerts.

External validation remains required before any broader claim. It includes coverage studies against documented Android behaviours, independent legal analysis of each mapping, usability work on notice comprehension and alert fatigue, accessibility review on target devices, security review of receipt handling, and field observation under an approved data-minimisation protocol. Until then the revised system is a reproducible, decoupled notification prototype: it makes a particular rule-owned combination visible for human review while explicitly leaving legal judgment and consequential intervention outside the App.
